\centerline{\bf \Huge Тренировочная олимпиада}
\centerline{\bf \Huge 8 класс}

\begin{flushright}
        \scriptsize{И помните: Империя заботится о вас!}
\end{flushright}

\problem{0}
\textit{Прикол от АА.}
Собрались как-то 157 учеников ЭлМШ, которым Расул Владимирович купил 100 пачек чокопаев(пачки были необычными двух видов: по 10 и 12 чокопаев в каждой соответсвенно). Сколько чокопаев съел каждый ученик, если не нашлось никаких двух учеников, которые съели разное количество чокопаев?

\problem{1}
На окне размером $1 \times 1$ сидит 99 мух. Испепеляющий лазер Ердема способен уничтожить круг радиуса $\dfrac{1}{9}$. Докажите, что за один выстрел Ердем способен сделать по мухам трипл килл.

\problem{2}
Докажите, что если сумма всех делителей одного числа равна сумме всех делителей другого числа и в то же время сумма чисел обратных делителям одного числа равна сумме чисел обратных делителям другого числа, то такие два числа равны.

\problem{3}
На координатной прямой нарисованы графики функций $y=3$ и $y=4$. На каждом из графиков друг под другом расположены 73 точки. Сколько существует прямоугольных треугольников в вершинах в этих точках?


\problem{4}
На клетчатом поле $7\times7$, возможно(а может и нет), стоит невидимая черная кошка, которая занимает квадрат $2\times 2$. Конфуций может расположить в некоторых клетках по чаше мудрости, а потом одновременно спросить у них, есть ли на их клетке чёрная кошка, после чего они одновременно ему отвечают. Какое наименьшее количество чаш нужно потратить, что гарантированно понять есть ли вообще на поле кошка и если да, то где она стоит?

\problem{5}
Квадратный трёхчлен $f(x)$ имеет ровно один корень. Кроме того, уравнение

$$
f(2 x-3)+f(3 x+1)=0
$$

имеет ровно один корень. Найдите корень трехчлена $f(x)$.